\documentclass[11pt,a4paper]{article}
\usepackage[utf8]{inputenc}
\usepackage[T1]{fontenc}
\usepackage[spanish]{babel}
\usepackage{amsmath,amssymb,mathtools}
\usepackage{geometry}
\usepackage{enumitem}
\usepackage{lmodern}
\usepackage{microtype}
\geometry{margin=2.4cm}
\setlength{\parindent}{0pt}
\setlength{\parskip}{0.45em}
\newcommand{\R}{\mathbb{R}}
\newcommand{\Ker}{\operatorname{Ker}}
\newcommand{\Ima}{\operatorname{Im}}
\newcommand{\rg}{\operatorname{rg}}
\newcommand{\tr}{\operatorname{tr}}
\newcommand{\Span}{\operatorname{L}}
\newcommand{\proy}{\operatorname{proy}}
\newcommand{\diag}{\operatorname{diag}}
\newcommand{\sg}{\operatorname{sg}}
\begin{document}

\begin{center}
{\LARGE\bfseries Álgebra Lineal}\\[0.4em]
{\Large Tema 6: El Espacio Vectorial Euclídeo}
\end{center}

\tableofcontents
\newpage

\section{Producto escalar}

Sea $V$ un espacio vectorial real. Un \textbf{producto escalar} es una forma bilineal:

\[
\langle\cdot,\cdot\rangle:V\times V\longrightarrow\R
\]

simétrica y definida positiva.

Para todo $u,v,w\in V$ y $\alpha\in\R$:

\[
\langle u,v\rangle=\langle v,u\rangle
\]

\[
\langle u+v,w\rangle
=
\langle u,w\rangle+\langle v,w\rangle
\]

\[
\langle \alpha u,v\rangle
=
\alpha\langle u,v\rangle
\]

\[
\boxed{
\langle u,u\rangle\geq0,
\qquad
\langle u,u\rangle=0\iff u=0.
}
\]

Un espacio vectorial real dotado de un producto escalar se denomina \textbf{espacio vectorial euclídeo}.

\subsection{Ejemplos}

En $\R^n$:

\[
\boxed{
\langle x,y\rangle
=
x_1y_1+\cdots+x_ny_n.
}
\]

En $\mathcal P_n(x)$:

\[
\boxed{
\langle p,q\rangle
=
\int_a^b p(x)q(x)\,dx.
}
\]

En $M_n(\R)$:

\[
\boxed{
\langle A,B\rangle
=
\tr(AB^t).
}
\]

\section{Matriz de Gram}

Sea:

\[
B=\{v_1,\ldots,v_n\}
\]

una base de $V$. La matriz del producto escalar en $B$ es:

\[
\boxed{
G_B=
\left(
\langle v_i,v_j\rangle
\right).
}
\]

Si:

\[
[x]_B=X,
\qquad
[y]_B=Y,
\]

entonces:

\[
\boxed{
\langle x,y\rangle
=
X^tG_BY.
}
\]

La matriz de Gram es simétrica y definida positiva.

\subsection{Cambio de base}

Si:

\[
X=PX',
\qquad
Y=PY',
\]

entonces:

\[
\boxed{
G_{B'}
=
P^tG_BP.
}
\]

\section{Norma y ángulo}

La norma de un vector es:

\[
\boxed{
\|v\|
=
\sqrt{\langle v,v\rangle}.
}
\]

Un vector es unitario si:

\[
\|v\|=1.
\]

Se verifican:

\[
\|\alpha v\|
=
|\alpha|\,\|v\|
\]

\[
\|u+v\|^2
=
\|u\|^2+\|v\|^2+2\langle u,v\rangle
\]

\[
\boxed{
\|u+v\|^2=\|u\|^2+\|v\|^2
\iff
\langle u,v\rangle=0
}
\]

\[
\boxed{
\|u+v\|^2+\|u-v\|^2
=
2\left(
\|u\|^2+\|v\|^2
\right)
}
\]

\[
\boxed{
|\langle u,v\rangle|
\leq
\|u\|\,\|v\|
}
\]

\[
\boxed{
\|u+v\|
\leq
\|u\|+\|v\|.
}
\]

El ángulo $\alpha$ entre dos vectores no nulos se define mediante:

\[
\boxed{
\cos\alpha
=
\frac{\langle u,v\rangle}
{\|u\|\,\|v\|}.
}
\]

\section{Ortogonalidad}

Dos vectores son ortogonales si:

\[
\boxed{
u\perp v
\iff
\langle u,v\rangle=0.
}
\]

Un sistema de vectores no nulos ortogonales dos a dos es linealmente independiente.

Para un subespacio $S\leq V$:

\[
\boxed{
S^\perp
=
\{v\in V:
\langle v,s\rangle=0
\quad
\forall s\in S
\}.
}
\]

Si:

\[
S=\Span\{u_1,\ldots,u_k\},
\]

entonces:

\[
\boxed{
S^\perp
=
\{v:
\langle v,u_i\rangle=0,
\ i=1,\ldots,k\}.
}
\]

Además:

\[
\boxed{
(S+T)^\perp
=
S^\perp\cap T^\perp.
}
\]

En dimensión finita:

\[
\boxed{
V=S\oplus S^\perp
}
\]

y:

\[
\boxed{
\dim S+\dim S^\perp=\dim V.
}
\]

\section{Bases ortogonales y ortonormales}

Una base:

\[
B=\{v_1,\ldots,v_n\}
\]

es ortogonal si:

\[
\langle v_i,v_j\rangle=0,
\qquad
i\neq j.
\]

Es ortonormal si además:

\[
\|v_i\|=1.
\]

Si $B$ es ortogonal, la matriz de Gram en $B$ es diagonal.

Si $B$ es ortonormal:

\[
\boxed{
G_B=I_n.
}
\]

A partir de una base ortogonal se obtiene una base ortonormal normalizando:

\[
\boxed{
e_i=
\frac{v_i}{\|v_i\|}.
}
\]

\section{Método de Gram--Schmidt}

Sea:

\[
B=\{v_1,\ldots,v_n\}.
\]

Definimos:

\[
u_1=v_1
\]

\[
u_2
=
v_2-
\frac{\langle v_2,u_1\rangle}
{\|u_1\|^2}u_1
\]

y, en general:

\[
\boxed{
u_i
=
v_i
-
\sum_{j=1}^{i-1}
\frac{\langle v_i,u_j\rangle}
{\|u_j\|^2}u_j.
}
\]

Los vectores:

\[
u_1,\ldots,u_n
\]

forman una base ortogonal y:

\[
\Span\{u_1,\ldots,u_i\}
=
\Span\{v_1,\ldots,v_i\}.
\]

Normalizando:

\[
\boxed{
e_i=
\frac{u_i}{\|u_i\|}
}
\]

se obtiene una base ortonormal.

\section{Proyecciones ortogonales}

\subsection{Proyección sobre un vector}

Para $u\neq0$:

\[
\boxed{
\proy_u(v)
=
\frac{\langle v,u\rangle}
{\|u\|^2}u.
}
\]

\subsection{Proyección sobre un subespacio}

Todo vector $v\in V$ se descompone de forma única como:

\[
\boxed{
v=s+w,
\qquad
s\in S,
\quad
w\in S^\perp.
}
\]

El vector $s$ es la proyección ortogonal de $v$ sobre $S$:

\[
s=\proy_S(v).
\]

Si:

\[
\{v_1,\ldots,v_k\}
\]

es una base ortogonal de $S$:

\[
\boxed{
\proy_S(v)
=
\sum_{i=1}^k
\frac{\langle v,v_i\rangle}
{\|v_i\|^2}v_i.
}
\]

La distancia de $v$ a $S$ es:

\[
\boxed{
d(v,S)
=
\|v-\proy_S(v)\|.
}
\]

\section{Matriz de una proyección ortogonal}

En $\R^n$ con el producto escalar usual, si las columnas de:

\[
A\in M_{n\times k}(\R)
\]

forman una base de $S$, la matriz de la proyección ortogonal es:

\[
\boxed{
P
=
A(A^tA)^{-1}A^t.
}
\]

Si las columnas de $A$ son ortonormales:

\[
\boxed{
P=AA^t.
}
\]

Toda matriz de proyección ortogonal cumple:

\[
\boxed{
P^2=P,
\qquad
P^t=P.
}
\]

Sus autovalores son:

\[
1
\]

sobre $S$ y:

\[
0
\]

sobre $S^\perp$.

\section{Método de mínimos cuadrados}

Consideramos un sistema sobredeterminado:

\[
AX=Y.
\]

El método de mínimos cuadrados busca el vector $X$ que minimiza:

\[
\|AX-Y\|^2.
\]

El residuo debe ser ortogonal al espacio columna de $A$:

\[
A^t(Y-AX)=0.
\]

Por tanto se obtienen las ecuaciones normales:

\[
\boxed{
A^tAX=A^tY.
}
\]

Si las columnas de $A$ son linealmente independientes:

\[
\boxed{
X=(A^tA)^{-1}A^tY.
}
\]

\subsection{Ajuste lineal}

Para ajustar:

\[
y=a+bx,
\]

se toma:

\[
A=
\begin{pmatrix}
1&x_1\\
\vdots&\vdots\\
1&x_m
\end{pmatrix},
\qquad
X=
\begin{pmatrix}
a\\b
\end{pmatrix},
\qquad
Y=
\begin{pmatrix}
y_1\\\vdots\\y_m
\end{pmatrix}.
\]

\subsection{Ajuste polinómico}

Para:

\[
y=a_0+a_1x+\cdots+a_nx^n,
\]

se utiliza la matriz:

\[
A=
\begin{pmatrix}
1&x_1&x_1^2&\cdots&x_1^n\\
\vdots&\vdots&\vdots&&\vdots\\
1&x_m&x_m^2&\cdots&x_m^n
\end{pmatrix}.
\]

\section{Endomorfismos simétricos y diagonalización ortogonal}

Sea $(V,\langle\cdot,\cdot\rangle)$ un espacio euclídeo.

Un endomorfismo $f:V\to V$ es simétrico si:

\[
\boxed{
\langle u,f(v)\rangle
=
\langle f(u),v\rangle
}
\]

para todos $u,v\in V$.

En una base ortonormal:

\[
\boxed{
f\text{ simétrico}
\iff
M_B(f)^t=M_B(f).
}
\]

Los autovalores de un endomorfismo simétrico son reales y los subespacios propios asociados a autovalores distintos son ortogonales.

\subsection{Teorema espectral}

Todo endomorfismo simétrico de un espacio euclídeo de dimensión finita admite una base ortonormal de autovectores.

Equivalentemente, toda matriz simétrica real es ortogonalmente diagonalizable:

\[
\boxed{
D=P^tAP=P^{-1}AP,
}
\]

con:

\[
\boxed{
P^tP=I.
}
\]

\section{Matrices y transformaciones ortogonales}

Una matriz es ortogonal si:

\[
\boxed{
P^tP=I.
}
\]

Equivalentemente:

\[
\boxed{
P^{-1}=P^t.
}
\]

Sus columnas forman una base ortonormal para el producto escalar usual.

Además:

\[
\boxed{
\det P=\pm1.
}
\]

El producto y la inversa de matrices ortogonales son ortogonales.

Una transformación ortogonal es una aplicación lineal que conserva el producto escalar:

\[
\boxed{
\langle f(x),f(y)\rangle
=
\langle x,y\rangle.
}
\]

Por tanto conserva también normas y ángulos.

Si $A=M_B(f)$ y $G_B$ es la matriz de Gram:

\[
\boxed{
G_B=A^tG_BA.
}
\]

\section{Matrices de Householder}

Sea $u\neq0$ y $U$ su vector columna.

La matriz de Householder asociada a $u$ es:

\[
\boxed{
H(u)
=
I-
2\frac{UU^t}{U^tU}.
}
\]

Se verifica:

\[
\boxed{
H^t=H,
\qquad
H^tH=I,
\qquad
\det H=-1.
}
\]

Representa una simetría respecto del hiperplano:

\[
u^\perp.
\]

Además:

\[
\boxed{
H(ku)=H(u),
\qquad
k\neq0.
}
\]

\section{Factorización QR}

Toda matriz con columnas linealmente independientes puede escribirse como:

\[
\boxed{
A=QR,
}
\]

donde:

\[
Q
\]

tiene columnas ortonormales y:

\[
R
\]

es triangular superior.

La factorización puede obtenerse aplicando Gram--Schmidt a las columnas de $A$ o mediante matrices de Householder.

\end{document}
